\documentclass[11pt]{article}
\usepackage[english,russian]{babel}
\usepackage[utf8]{inputenc}
\usepackage[T2A]{fontenc}
\usepackage{url}
\usepackage{graphicx,DCCN2026_ru}

\pagestyle{fancy}
\fancyhead{}
\fancyfoot{}

\linespread{1.0}

\usepackage{amsmath,amssymb}
\usepackage{booktabs}

\makeatletter
\fancyhead[RO]{\small DCCN 2026\\ {21--25 сентября 2026}}
\fancyhead[LO]{\small Без авторов (double-blind)\\
Иерархическая STRL-архитектура управления}

\c@page=1
\makeatother


\title{Иерархическая STRL-архитектура управления для мобильного робота:
интеграция TD3, LQR-residual и обхода препятствий с использованием карты}

\author[1]{\small }
\author[2]{\small }

\email{ }

\begin{document}

\udc{004.896}

{\let\newpage\relax\maketitle}
\vskip -1.5em
%\footnotetext{Работа выполнена при поддержке ...}


\begin{abstract}
Рассматривается задача навигации мобильного робота с дифференциальным
приводом в плоской среде с препятствиями и переменными свойствами
поверхности. Предложена реализация трёхуровневой иерархической архитектуры,
соотносимой с парадигмой Strategic-Tactical-Reactive Levels (STRL) школы
Осипова-Панова и объединяющей реактивную политику TD3, тактический слой
LQR-residual и супервизорный модуль обхода препятствий на основе известной
карты. Проведён ablation-анализ на десяти сидах для двух семейств сценариев:
семи конфигураций препятствий и трёх типов поверхностей. Показано, что
распределение функций между уровнями архитектуры существенно зависит
от структуры среды: супервизорный уровень разгружает реактивный в задачах
с известной геометрией препятствий, но не компенсирует физических
ограничений нижнего уровня в условиях слабого сцепления колёс. Все
результаты воспроизводимы со 100\,\% детерминированностью.

\keywords{обучение с подкреплением, TD3, LQR, иерархическое управление,
мобильная робототехника, STRL, residual policy, когнитивный агент}
\end{abstract}

\section{Введение}
\label{sec:intro}

Автономная навигация мобильного робота в среде с препятствиями и переменными
свойствами поверхности~--- задача, в которой соприкасаются несколько направлений
современных исследований искусственного интеллекта: обучение с подкреплением,
классическое управление, иерархическое планирование. Каждое из направлений
решает её на своём уровне абстракции и в своих допущениях, и вопрос об
их интеграции в единую архитектуру управления остаётся открытым.

В традиции отечественной школы когнитивных агентов (школа Осипова-Панова,
ФИЦ ИУ РАН) указанный вопрос рассматривается через парадигму STRL~---
Strategic-Tactical-Reactive Levels~\cite{kiselev2020dissertation,gat1998three}.
В STRL поведение агента распределено по трём уровням, действующим в существенно
разных временных масштабах: реактивный уровень управляет моторами на
миллисекундной шкале, тактический~--- координирует локальное движение
на шкале секунд, стратегический~--- планирует поведение на длинных горизонтах
на основе модели мира. В диссертационной работе~\cite{kiselev2020dissertation}
STRL-парадигма была реализована для мобильного робота МП-РМ <<Зарница>>
с Q-обучением на реактивном уровне; продолжение этой линии в работе
\cite{veitsenfeld2024} расширяет внимание на мультимодальные сигналы
и переходит к VLM-моделям как кандидатам на роль более интегрированного
агента.

В настоящей работе предложена реализация STRL-парадигмы для случая
непрерывного управления мобильным роботом с дифференциальным приводом.
Реактивный уровень реализован обученной методом TD3~\cite{fujimoto2018td3}
нейросетевой политикой; тактический~--- LQR-регулятором по линеаризованной
модели бокового движения; супервизорный (в нашей реализации играющий роль
стратегического уровня, см. раздел~\ref{sec:supervisory})~--- модулем
обхода препятствий на основе известной карты. Архитектура
проверена в физическом симуляторе PyBullet на платформе Husky на двух
семействах сценариев: семь конфигураций препятствий с разной геометрией
и три типа поверхности.

\textbf{Вклад работы.} (1)~Реализация трёхуровневой иерархической архитектуры,
соотносимой со STRL-парадигмой, для непрерывного управления мобильным роботом
в среде с препятствиями и переменными свойствами поверхности. (2)~Эмпирический
анализ распределения функций между уровнями на двух семействах сценариев
($7 \times 5 + 6 \times 5 = 65$ ablation-ячеек, $n = 10$ запусков в каждой,
100\,\% воспроизводимость). Анализ показывает, что супервизорный уровень
с моделью среды разгружает реактивный в задачах с известной геометрией
препятствий, но не компенсирует физические ограничения нижнего уровня
в условиях слабого сцепления колёс с поверхностью. (3)~Открытый репозиторий
с воспроизводимой реализацией всех описанных уровней, обученной моделью
и протоколом экспериментов.

Дальнейший текст организован следующим образом. В разделе~\ref{sec:related}
приводится обзор работ по основным используемым методам и парадигмам.
В разделе~\ref{sec:method} формулируется задача и описывается архитектура
управления. Раздел~\ref{sec:experiments} содержит результаты ablation-анализа
на двух семействах сценариев. Раздел~\ref{sec:discussion} обсуждает
распределение функций между уровнями архитектуры и ограничения подхода.
Раздел~\ref{sec:conclusion} формулирует выводы и направления дальнейших
исследований.


\section{Обзор работ}
\label{sec:related}

\subsection{Обучение с подкреплением для непрерывного управления}
\label{sec:rw_rl}

Задача управления мобильным роботом в непрерывном пространстве действий
относится к классу проблем, для которых табличные алгоритмы динамического
программирования и DQN~\cite{mnih2015dqn} неприменимы из-за невозможности
представить $Q(s, a)$ для бесконечного множества действий. В таких задачах
используют параметрические аппроксиматоры $Q_\phi$ и $\pi_\theta$, обучаемые
методами актёр-критика: DDPG, TD3~\cite{fujimoto2018td3}, SAC, PPO. В настоящей
работе используется TD3 как off-policy метод с подтверждённой устойчивостью
обучения на робототехнических задачах и доступной эталонной реализацией
в \verb|stable-baselines3|~\cite{raffin2021sb3}.

\subsection{Residual Policy и LQR в управлении мобильными роботами}
\label{sec:rw_residual}

Идея residual policy learning~\cite{silver2018residual,johannink2019residual}
возникла как ответ на ограничение обучения с подкреплением с нуля и
классических контроллеров: первое требует на порядки больше взаимодействий
со средой, второе теряет качество при контактах, трении и неучтённой
нелинейности. Residual policy комбинирует обучаемую политику с
классическим контроллером, опираясь на сильные стороны обоих. Линейно-квадратичный
регулятор (LQR) применяется к нелинейной кинематике дифференциального привода
через линеаризацию в координатах ошибки относительно опорной траектории.
В настоящей работе реализован lateral LQR с упрощением $e_y = 0$,
оправданным постановкой навигации к точечной цели.

\subsection{Иерархический RL и waypoint-планирование}
\label{sec:rw_hrl}

Идея разбиения задачи управления на уровни различной временной абстракции
восходит к option framework~\cite{sutton1999options}. Современные
иерархические схемы reinforcement learning представлены HIRO~\cite{nachum2018hiro},
в котором цели задаются в исходном пространстве с off-policy correction,
и Director~\cite{hafner2022director}, в котором цели заданы в латентном
пространстве world model. Между двумя крайностями располагается семейство
схем с фиксированной верхнеуровневой структурой~--- waypoint-планировщики
и классические геометрические алгоритмы навигации. В настоящей работе
супервизорный уровень реализован простейшим геометрическим способом,
что позволяет изолировать вклад модельного знания, отделяя его от вклада
сложности алгоритма планирования.


\subsection{STRL-парадигма и школа Осипова-Панова}
\label{sec:rw_strl}

Идея многоуровневой архитектуры мобильного робота сложилась в работах
по поведенческой робототехнике 1990-х годов: Brooks обосновал subsumption
architecture, Gat~\cite{gat1998three} систематизировал трёхуровневую
организацию (deliberator-sequencer-controller), Arkin предложил
поведенческие конечные автоматы для координации реактивных модулей.
В школе Осипова-Панова эта линия получила развитие в направлении
когнитивных агентов со знаковой картиной мира: операции на знаковых
сетях с кортежем $s = \langle n, p, m, a \rangle$ (имя, образ, значение,
смысл)~\cite{osipov2018sign} соединяют сенсомоторный опыт с понятийной
структурой картины мира. Диссертационная работа~\cite{kiselev2020dissertation}
реализует Strategic-Tactical-Reactive Levels (STRL) для робототехнической
платформы МП-РМ, с Q-обучением~\cite{kiselev2020} атомарных действий
на реактивном уровне и иерархическим планированием на верхних. Прямой
предшественник настоящей работы~--- статья~\cite{veitsenfeld2024}~---
расширяет программу школы в сторону мультимодальных сигналов и
visual-language моделей; в её заключительной части обозначено направление
дополнения VLM компонентой обучения с подкреплением и переход к интерактивной
среде. Настоящая работа развивает указанное направление, реализуя
парадигму STRL для непрерывного управления мобильным роботом в физическом
симуляторе, и проводит количественный анализ распределения функций между
уровнями~--- вопрос, который в дискретной постановке~\cite{kiselev2020dissertation}
не мог быть исследован систематически из-за ограничений тестовой платформы.



\section{Метод}
\label{sec:method}

\subsection{Постановка задачи}
\label{sec:problem}

Рассматривается задача навигации мобильного робота к случайной целевой точке
на плоской арене, возможно содержащей препятствия известной геометрии и
варьируемой поверхностью под колёсами. Среда моделируется в физическом
симуляторе PyBullet~\cite{coumans2021pybullet}; в качестве платформы
используется Husky~--- четырёхколёсный робот с дифференциальным приводом,
входящий в стандартную поставку симулятора. Управление осуществляется на
уровне целевых угловых скоростей колёс, причём левая и правая пары
рассматриваются как два управляющих канала, что эквивалентно стандартной
формулировке differential-drive.

Задача формулируется как процесс принятия решений в марковской модели
$\langle S, A, T, R, \gamma \rangle$. \textbf{Состояние} $s \in S$ образует
вектор $o \in \mathbb{R}^{n_o}$ позы и кинематики робота, дополненный
относительным положением цели и (для среды с препятствиями) лидарными
измерениями; конкретный состав вектора зависит от экспериментальной
установки. \textbf{Действие} $a \in [-1, 1]^2$~--- нормированные целевые
скорости левой и правой колёсных пар, преобразуемые в физические
$\omega_{\mathrm{wheel}} \in [-\omega_{\max}, \omega_{\max}]$ с
$\omega_{\max} = 10$~рад/с (линейная скорость робота до 1.7~м/с).
\textbf{Функция перехода}~$T$ определяется физикой PyBullet с шагом
симуляции $1/240$~с; шаг управления $1/20$~с, то есть одно действие
удерживается двенадцать физических шагов. \textbf{Дисконт}~$\gamma = 0.99$.

\textbf{Арена}~--- горизонтальная плоскость размером $10 \times 10$~м; робот
стартует из центра, цель сэмплируется на расстоянии от 2 до 5~м от старта
в случайном направлении, $d_{\mathrm{goal}} = 0.4$~м~--- радиус достижения
цели. \textbf{Условия завершения эпизода}: достижение цели (success),
переворот робота на $> 60^{\circ}$ по тангажу или крену (failure), выход за
пределы арены (failure), столкновение с препятствием (failure) или
исчерпание лимита 500 шагов управления, что соответствует примерно 25~с
реального времени (timeout).

\textbf{Функция вознаграждения} строится плотной~--- сумма штрафа за
расстояние до цели и постоянного штрафа за каждый шаг, с дополнительными
терминальными членами и средоспецифическими компонентами. При наличии
препятствий к базовой части добавляются плотный штраф за близость
к препятствиям и терминальный штраф за столкновение.

Постановка отличается от стандартных бенчмарков непрерывного управления
типа Reacher или HalfCheetah тем, что робот наделён неголономной
кинематикой дифференциального привода и решает задачу пространственной
навигации, а не локальной манипуляции. Это сближает её с практическими
задачами мобильной робототехники.


\subsection{Иерархическая архитектура управления}
\label{sec:architecture}

В качестве организующего принципа управления используется парадигма
многоуровневой архитектуры, восходящая к ранним работам по
поведенческой робототехнике и реактивным архитектурам~\cite{gat1998three},
и развитая для случая когнитивных агентов в школе Осипова-Панова и
конкретно в диссертационной работе~\cite{kiselev2020dissertation},
где обозначена как Strategic--Tactical--Reactive Levels (STRL).
В STRL-парадигме поведение агента распределено по трём уровням,
действующим в существенно разных временных масштабах и оперирующим
разными моделями среды.

В настоящей работе реализована адаптация STRL для непрерывного управления
мобильным роботом с дифференциальным приводом (рисунок~\ref{fig:strl_arch}).
Три уровня выделяются следующим образом.

\textbf{Реактивный уровень} реализован обученной методом TD3~\cite{fujimoto2018td3}
нейросетевой политикой $\pi_{\mathrm{TD3}}: o \mapsto a$, отображающей текущее
наблюдение в нормированное двухкомпонентное действие. Этот уровень действует
на каждом шаге управления (период 50~мс), оперирует только текущим наблюдением
и в одиночку способен решать задачу навигации с приемлемым качеством
в простых средах.

\textbf{Тактический уровень} реализован LQR-регулятором с фиксированным
коэффициентом усиления, построенным на линеаризованной модели бокового
движения differential-drive. Контроллер возвращает скалярную коррекцию
$u_{\mathrm{LQR}} \in [-u_{\max}, u_{\max}]$ к угловой скорости робота,
комбинируемую с действием реактивного уровня в схеме Residual
Policy~\cite{silver2018residual,johannink2019residual}. Тактический уровень
действует на том же временном масштабе, что и реактивный, но опирается
на явную аналитическую модель кинематики, что обеспечивает гладкость
рулевого управления независимо от качества обучения нижнего уровня.

\textbf{Супервизорный уровень} реализован модулем обхода препятствий
с использованием известной карты их положений (MapAware). Этот уровень
получает на вход геометрию препятствий и вычисляет необходимость
вмешательства: при приближении к любому препятствию в передней
полусфере на расстояние ближе $d_{\mathrm{takeover}} = 1.5$~м управление
полностью передаётся супервизорному контроллеру до выхода из радиуса
$d_{\mathrm{release}} = 2.5$~м. Супервизорный уровень действует
эпизодически~--- по необходимости перехватывая управление в опасных
конфигурациях, тогда как в безопасной геометрии полностью отдаёт
контроль реактивному и тактическому уровням. Содержательно этот уровень
играет роль того, что в классической STRL называется стратегическим
уровнем, с тем уточнением, что в нашей реализации он не планирует
последовательность подцелей, а supervisorно вмешивается на основе
модели среды; в дальнейшем мы обсуждаем эту особенность в разделе
ограничений.

Финальное действие, передаваемое в среду на каждом шаге управления,
определяется правилом приоритета: при активном супервизорном уровне
$a = a_{\mathrm{MapAware}}$; иначе
\begin{equation}
\label{eq:combine}
a = a_{\mathrm{TD3}} + \alpha \cdot \begin{bmatrix} -1 \\ +1 \end{bmatrix} u_{\mathrm{LQR}},
\end{equation}
где знаки $-1, +1$ преобразуют скалярную поправку $u_{\mathrm{LQR}}$
в дифференциальный сигнал на колёса, а $\alpha \in [0, 1]$~---
фиксированный весовой коэффициент. Результат покомпонентно ограничивается
отрезком $[-1, 1]$. При $\alpha = 0$ и неактивном супервизоре архитектура
вырождается в чистый TD3. В работе используется $\alpha = 0.2$, выбранное
эмпирически как компромисс между подавлением высокочастотных колебаний
рулевого управления и сохранением стратегических решений политики.

% Рисунок 1: schema strl_arch — TODO: создать
\begin{figure}[h]
   \centering
   \fbox{\parbox{0.8\textwidth}{\centering\vspace{2cm} [Рисунок 1: схема STRL-архитектуры]\\ Reactive (TD3) / Tactical (LQR) / Supervisory (MapAware) \\ \vspace{2cm}}}
   \caption{Схема трёхуровневой иерархической архитектуры управления.
   Реактивный уровень (TD3) действует на каждом шаге; тактический (LQR-residual)
   корректирует курсовое движение; супервизорный (MapAware) перехватывает
   управление при приближении к препятствиям, положения которых известны.}
\label{fig:strl_arch}
\end{figure}

\subsection{Реактивный уровень: TD3}
\label{sec:reactive}

В качестве базовой политики реактивного уровня используется
алгоритм TD3~\cite{fujimoto2018td3}~--- off-policy актёр-критик
для непрерывного управления с двойным критиком и задержанным обновлением
актёра. Актёр и оба критика реализованы как полносвязные нейронные сети
со скрытыми слоями размерности 400 и 300 и нелинейностью ReLU; выход
актёра пропускается через $\tanh$. Архитектура воспроизводит стандартную
конфигурацию TD3 в реализации Stable-Baselines3~\cite{raffin2021sb3}.

Гиперпараметры обучения: batch~256, оптимизатор Adam, скорость обучения
$10^{-3}$, $\tau = 0.005$, задержка обновления актёра $d = 2$, target-smoothing
$\sigma = 0.2$ с обрезкой $\pm 0.5$. Шум исследования~--- гауссов
$\mathcal{N}(0, 0.3)$. Размер replay-буфера $10^5$ в пустой арене и
$2\cdot 10^5$ в среде с препятствиями. На каждый шаг сбора опыта приходится
один градиентный шаг для каждого критика и (с учётом задержки) один шаг
для актёра.

Наблюдение реактивного уровня включает позу робота, его линейные и угловые
скорости, ориентацию (косинус и синус угла рыскания) и относительное
положение цели. Такое наблюдение является \emph{привилегированным}~---
оно эквивалентно идеальной локализации робота и точному знанию положения
цели~--- что отделяет вопросы архитектуры управления от вопросов восприятия.
В среде с препятствиями вектор расширяется лидарными измерениями
(16 лучей в передней полусфере). Каждая модель обучается в одной заданной
среде до выхода метрики evaluation reward на устойчивое плато; между
моделями для разных сред веса не переносятся.


\subsection{Тактический уровень: LQR-residual}
\label{sec:tactical}

Тактический уровень построен на линеаризованной модели бокового движения
differential-drive. В системе координат, связанной с прямой, ведущей от
робота к текущей цели, состояние ошибки выбирается двумерным:
$e = [e_y, e_\theta]^\top$, где $e_y$~--- поперечное отклонение от опорной
прямой, $e_\theta$~--- отклонение курса. При фиксированной опорной линейной
скорости $v_{\mathrm{ref}}$ и малых отклонениях линеаризованная модель
ошибки принимает вид
\begin{equation}
\label{eq:lqr_model}
\dot{e} = \begin{bmatrix} 0 & v_{\mathrm{ref}} \\ 0 & 0 \end{bmatrix} e
         + \begin{bmatrix} 0 \\ -1 \end{bmatrix} \omega.
\end{equation}

LQR находит оптимальное управление $\omega = -K e$, минимизирующее
квадратичный функционал
\begin{equation}
\label{eq:lqr_cost}
J = \int_0^\infty \left( e^\top Q e + \omega^\top R \omega \right) dt,
\end{equation}
с весовыми матрицами $Q = \mathrm{diag}(2.0,\ 1.0)$ и $R = 5.0$. Соотношение
весов выбрано так, чтобы штраф на управление превышал штраф на отклонение
курса: контроллер выдаёт мягкие, не доминирующие коррекции. Решение
алгебраического уравнения Риккати выполнено численно средствами
\verb|scipy.linalg.solve_continuous_are|; матрица $K$ вычисляется один раз
при инициализации.

В реализации сделано существенное упрощение: компонента $e_y$ обнуляется
на каждом шаге управления. Причина~--- постановка задачи навигации
к точечной цели: опорная прямая переопределяется на каждом шаге как линия
из текущей позы робота в текущую цель, и накопленного поперечного
отклонения относительно этой динамической прямой не возникает. Закон
управления вырождается в $\omega = -k_\theta\, e_\theta$, где $k_\theta$~---
соответствующий элемент матрицы $K$. Структурно это P-регулятор по углу
с коэффициентом, выведенным из LQR-задачи. Такое сознательное упрощение
оправдано тем, что в схеме Residual Policy LQR играет вспомогательную
роль и должен подавлять высокочастотные колебания рулевого управления,
не вмешиваясь в стратегические решения политики.

Скалярный выход $\omega$ ограничивается отрезком $[-u_{\max}, u_{\max}]$
с $u_{\max} = 0.3$ для гарантии того, что при патологических значениях
$e_\theta$ поправка остаётся в безопасном диапазоне.


\subsection{Супервизорный уровень: обход с использованием карты}
\label{sec:supervisory}

Супервизорный уровень получает на вход карту препятствий~---
список троек $\{(o^x_i, o^y_i, r_i)\}$ цилиндрических препятствий с
положениями центров и радиусами. В каждый момент времени уровень
вычисляет расстояния и углы до всех препятствий в передней полусфере
относительно курса робота, и применяет двупороговую логику:
управление перехватывается, если минимальное расстояние до края
препятствия в полусфере $\pm \pi/2$ относительно курса оказывается
ниже $d_{\mathrm{takeover}} = 1.5$~м, и возвращается реактивно-тактической
паре, когда минимальное расстояние превышает $d_{\mathrm{release}} = 2.5$~м.
Гистерезис между порогами предотвращает дребезг управления вблизи границы.

Во время активного супервизорного режима контроллер задаёт прямолинейное
движение со скоростью $v_{\mathrm{sup}} = 0.4$~м/с и угловую скорость,
поворачивающую робот от ближайшего препятствия:
\begin{equation}
\label{eq:supervisory}
\omega_{\mathrm{sup}} = k_{\mathrm{turn}} \cdot \mathrm{sign}(-\theta_{\mathrm{nearest}}),
\end{equation}
где $\theta_{\mathrm{nearest}}$~--- локальный угол на ближайшее препятствие
в передней полусфере, $k_{\mathrm{turn}} = 1.5$. Логика преднамеренно
проста: в задачах, где геометрия препятствий известна, эффективное
поведение часто достижимо без оптимизационного планировщика. Это позволяет
изолировать вклад \emph{знания карты} как такового, отделяя его от вклада
сложности конкретного алгоритма планирования.

Содержательно описанный супервизорный уровень не реализует стратегический
синтез плана в классической STRL-парадигме~\cite{kiselev2020dissertation}:
он не предлагает последовательность подцелей, а перехватывает управление
в опасных конфигурациях. Это сознательное упрощение, преследующее цель
выделить роль модельного знания (известной карты) как такового. Расширение
до полноценного стратегического уровня~--- адаптивного планировщика
на основе известной карты~--- естественное направление развития работы,
обсуждаемое в разделе~\ref{sec:conclusion}.

Альтернативой супервизорному режиму на основе карты в работе также
рассмотрена реактивная heuristic-схема обхода по лидарным дальностям
(LidarAvoid). Сравнение этих двух режимов друг с другом и с архитектурами
без супервизора составляет основной экспериментальный материал
раздела~\ref{sec:experiments}.
\section{Эксперименты}
\label{sec:experiments}

\subsection{Конфигурация и метрики}
\label{sec:exp_setup}

Все эксперименты проведены на вычислительном сервере (8 ядер Ryzen 9 9950X,
16 ГБ RAM, без GPU; обучение TD3 на CPU средствами Stable-Baselines3).
Симулятор PyBullet 3.2 запускался в безголовом режиме с параллельным
исполнением шести независимых сред через \verb|SubprocVecEnv|. Все
эксперименты с препятствиями детерминированы по сидам $\{200, 201, \ldots, 209\}$;
для каждой конфигурации архитектуры и сценария проведено по $n = 10$
запусков, и таблицы ниже сообщают долю успешных эпизодов из десяти.

Основной метрикой служит \emph{success rate}~--- доля эпизодов, завершившихся
достижением целевой точки до timeout'а 500 шагов и без столкновения
с препятствием. Дополнительно для качественного анализа фиксируются:
средняя длина эпизода (число шагов управления), мера гладкости управления
(средний модуль разности между последовательными командами на колёса),
максимальная развитая линейная скорость, доля эпизода в супервизорном
режиме.

Сравниваются пять архитектурных режимов, соответствующих включению/выключению
тактического и супервизорного уровней:
\begin{itemize}
\item \textbf{TD3}~--- только реактивный уровень.
\item \textbf{TD3+LQR}~--- реактивный и тактический уровни.
\item \textbf{TD3+LidarAvoid}~--- реактивный и реактивный супервизор по лидару.
\item \textbf{TD3+MapAvoid}~--- реактивный и модельный супервизор по карте.
\item \textbf{TD3+MapAvoid+LQR}~--- все три уровня.
\end{itemize}

\subsection{Эксперимент A: множественные конфигурации препятствий}
\label{sec:exp_a}

В первой серии экспериментов оценивается зависимость распределения функций
от геометрии препятствий. Используется семь детерминированных конфигураций
препятствий, представляющих качественно разные ситуации навигации:
одиночное препятствие на прямой к цели; два препятствия со смещением;
коридор из трёх препятствий; барьер с проходом; слалом из четырёх
препятствий; диагональное расположение цели; широкое препятствие.
В каждой конфигурации цель располагается в фиксированной точке,
геометрия препятствий полностью известна супервизорному уровню.

\begin{table}[h]
\begin{center}
\begin{tabular}{lccccc}
\toprule
Сценарий & TD3 & +LQR & +Lidar & +Map & +Map+LQR \\
\midrule
single\_on\_line     & 0/10 & 0/10 & 0/10 & \textbf{10/10} & \textbf{10/10} \\
two\_offset         & 0/10 & 0/10 & 0/10 & \textbf{10/10} & \textbf{10/10} \\
three\_corridor     & 0/10 & 0/10 & 0/10 & 0/10           & 0/10           \\
barrier\_with\_gap  & \textbf{10/10} & \textbf{10/10} & \textbf{10/10} & 0/10 & 0/10 \\
slalom              & 0/10 & 0/10 & 0/10 & 0/10           & 0/10           \\
diagonal\_goal      & 10/10 & 10/10 & 10/10 & 10/10        & 10/10         \\
wide\_obstacle      & 0/10 & 0/10 & \textbf{10/10} & 10/10 & 10/10         \\
\bottomrule
\end{tabular}
\caption{Эксперимент A: success rate (из 10 запусков на сидах 200--209)
для семи конфигураций препятствий и пяти архитектурных режимов.
Все ячейки 100\% детерминированы: повторные запуски в идентичной
конфигурации дают идентичный исход.}
\label{tab:exp_a}
\end{center}
\end{table}

Данные таблицы~\ref{tab:exp_a} демонстрируют существенно неравномерное
распределение успеха по сценариям и режимам. Сценарий \emph{diagonal\_goal}
не содержит препятствий на прямой к цели, и все пять режимов справляются
с ним полностью. Это служит саниксом: реактивный уровень не деградирует
при включении тактического или супервизорного уровней. Сценарии
\emph{three\_corridor} и \emph{slalom} не решаются ни одним режимом~---
они выходят за пределы возможностей данной архитектуры и фиксируют верхнюю
границу её применимости.

Содержательное различие проявляется на оставшихся четырёх сценариях.
В \emph{single\_on\_line} и \emph{two\_offset} разница между режимами
максимальна: только включение супервизорного уровня (TD3+MapAvoid,
TD3+MapAvoid+LQR) даёт 100\% успех, в то время как реактивный и тактический
уровни в одиночку или в комбинации не справляются. На уровне траекторий
видно, что обученная TD3-политика идёт по прямой в финальную цель,
не учитывая препятствие, и сталкивается с ним; LQR-residual подавляет
колебания курса, но не меняет общую стратегию. Только супервизорный
уровень с явным знанием карты способен перенаправить движение, не нарушая
итоговый успех.

Обратный паттерн наблюдается в сценарии \emph{barrier\_with\_gap}, где
прямой путь к цели проходит через узкий проход в барьере. Здесь все
архитектуры, опирающиеся на супервизорный уровень MapAvoid, не справляются
(0/10), а полностью реактивные (TD3, TD3+LQR) и реактивный
с лидарным супервизором (TD3+Lidar) справляются полностью. Причина видна
в логике MapAvoid: широкий радиус takeover ($d_{\mathrm{takeover}} = 1.5$~м)
заставляет супервизор инициировать обход барьера в обход прохода,
тогда как геометрический проход в барьере проходим напрямую. Реактивная
политика, обученная на разнообразной геометрии, проходит через проход
успешно.

Наконец, в сценарии \emph{wide\_obstacle} наибольшую устойчивость показывает
лидарный супервизор: широкое препятствие не маркировано в карте как
протяжённое (используется аппроксимация цилиндром малого радиуса),
и MapAvoid не активируется на нужной дистанции, тогда как лидарные лучи
фиксируют препятствие напрямую. Это указывает на ограниченность модельного
знания: оно полезно, но требует адекватной геометрической репрезентации
препятствий.

\textbf{Воспроизводимость.} Все сорок ячеек в таблице~\ref{tab:exp_a}
проверены повторными запусками в идентичной конфигурации. Результаты
полностью детерминированы: каждая ячейка имеет однозначный исход,
$10/10$ или $0/10$, без промежуточных значений на десяти сидах.
Это свойство сцены, а не статистики, и оно следует из детерминированности
PyBullet при фиксированных стартовых условиях, фиксированных весах
обученной политики и отсутствия стохастичности в супервизорном и
тактическом уровнях.


\subsection{Эксперимент B: переменные свойства поверхности}
\label{sec:exp_b}

Во второй серии оценивается влияние свойств опорной поверхности на
работу архитектуры. Реализованы три условия: \emph{NORMAL} (коэффициент
трения 0.9, базовая поверхность), \emph{ICE} (коэффициент трения 0.05,
скользкая поверхность, имитирующая лёд) и \emph{SAND} (коэффициент трения
0.6 с дополнительным сопротивлением качению, имитирующая мягкий грунт).
Каждое условие тестируется в двух конфигурациях: \emph{empty}~--- пустая
арена с целью, \emph{with\_obstacle}~--- арена с препятствием, для которого
известна карта. Каждая ячейка содержит десять запусков на сидах 200--209.

\begin{table}[h]
\begin{center}
\begin{tabular}{llccccc}
\toprule
Поверхность & Конфиг. & TD3 & +LQR & +Lidar & +Map & +Map+LQR \\
\midrule
NORMAL & empty     & 10/10 & 10/10 & 10/10 & 10/10 & 10/10 \\
NORMAL & + obst.   & 0/10  & 0/10  & 0/10  & \textbf{10/10} & \textbf{10/10} \\
ICE    & empty     & 10/10 & 10/10 & 10/10 & 10/10 & 10/10 \\
ICE    & + obst.   & 0/10  & 0/10  & 0/10  & 0/10  & 0/10  \\
SAND   & empty     & 10/10 & 10/10 & 10/10 & 10/10 & 10/10 \\
SAND   & + obst.   & 0/10  & 0/10  & 0/10  & 10/10 & 10/10 \\
\bottomrule
\end{tabular}
\caption{Эксперимент B: success rate (из 10 запусков) для трёх типов
поверхности и двух конфигураций. Модель TD3 v3 обучена на детерминированном
сценарии с препятствием на NORMAL и используется во всех ячейках без
переобучения.}
\label{tab:exp_b}
\end{center}
\end{table}

Таблица~\ref{tab:exp_b} разделяет два эффекта. В пустой арене реактивный
уровень одинаково успешен на всех трёх поверхностях, хотя максимальная
развитая скорость различается на порядок (на ICE робот вынужден двигаться
вдвое медленнее, чтобы не соскользнуть; на SAND~--- медленнее на четверть).
Это означает, что обученная политика обладает достаточной общностью для
адаптации линейной скорости под трение, хотя обучение проводилось только
на NORMAL.

В присутствии препятствия картина существенно меняется. На NORMAL и SAND
включение супервизорного уровня MapAvoid полностью решает задачу
(10/10 в обоих случаях), что согласуется с результатами
эксперимента~A в сценарии \emph{single\_on\_line}. Однако на ICE
с препятствием ни один из режимов не решает задачу: при попытке обхода
препятствия низкое трение приводит к проскальзыванию колёс, робот
заносит, и столкновение происходит даже при идеальной стратегии обхода.
Это указывает на принципиальное ограничение архитектуры: \emph{знание
карты не компенсирует физическое ограничение нижнего уровня}.

Для проверки этой гипотезы было проведено дополнительное дообучение
TD3-политики на расширенной выборке поверхностей (continued training
на 700~тыс. шагов с распределением 50\%/30\%/20\% между NORMAL/SAND/ICE).
Получившаяся модель показала улучшение ICE+obstacle с 0/10 до 5/10
в комбинации с тактическим уровнем, но регрессию на NORMAL+obstacle
с 10/10 до 7/10 в режиме MapAvoid+LQR. Это эмпирически фиксирует
trade-off между общностью и пиковой скоростью на отдельно взятой
поверхности, известный в литературе для transfer learning, и формирует
естественное ограничение применимости архитектуры в текущем виде.
\section{Обсуждение}
\label{sec:discussion}

Эмпирические результаты раздела~\ref{sec:experiments} позволяют сформулировать
несколько содержательных наблюдений о распределении функций между уровнями
рассмотренной иерархической архитектуры.

\textbf{Реактивный уровень компетентен в условиях свободного движения.}
В пустой арене и на всех трёх типах поверхности (NORMAL, ICE, SAND)
обученная TD3-политика достигает 100\,\% успеха, демонстрируя способность
адаптировать линейную скорость под условия трения без переобучения.
Это совпадает с типичной областью применения reactive control в STRL-парадигме:
действие на коротком горизонте без обращения к модели среды.

\textbf{Тактический уровень не меняет стратегические решения.} Включение
LQR-residual на верхушке реактивного уровня не меняет success rate ни в одной
из исследованных ячеек~--- ни в эксперименте~A (таблица~\ref{tab:exp_a}),
ни в эксперименте~B (таблица~\ref{tab:exp_b}). Качественно LQR подавляет
высокочастотные колебания курсового управления, что фиксируется в метрике
гладкости (снижение на $\sim 10\,\%$, согласно дополнительным измерениям),
но не влияет на способность преодолевать препятствия или соскальзывать
с поверхности. Это согласуется с теоретической ролью тактического уровня
в STRL~\cite{kiselev2020dissertation} и подтверждает, что упрощение
$e_y = 0$ в лateral LQR не лишает регулятор полезной функции.

\textbf{Супервизорный уровень разгружает реактивный там, где
модельное знание адекватно.} Эксперимент~A показывает, что в сценариях,
где препятствие лежит на прямой к цели (\emph{single\_on\_line},
\emph{two\_offset}), реактивный уровень не справляется в одиночку,
но включение супервизорного с MapAware-логикой обеспечивает 100\,\% успех.
Это типичная демонстрация полезности стратегического уровня в STRL: знание
среды позволяет принимать решения, которых реактивный уровень,
действующий по сенсорному наблюдению, принять не может. Однако
в сценарии \emph{barrier\_with\_gap} мы наблюдаем обратный эффект:
супервизор перехватывает управление при приближении к барьеру и
заставляет робот обойти его, тогда как геометрически проход проходим.
Это указывает на чувствительность супервизорного уровня к параметрам
геометрической модели препятствий и к структуре алгоритма обхода.

\textbf{Модельное знание не компенсирует физическое ограничение нижнего
уровня.} Эксперимент~B даёт принципиальное наблюдение: на ICE с препятствием
ни один из режимов не решает задачу, включая полностью трёхуровневую
архитектуру TD3+MapAvoid+LQR. Супервизор предлагает корректную стратегию
обхода, тактический уровень обеспечивает гладкое управление, но низкое
трение приводит к проскальзыванию и столкновению независимо от стратегии.
Иными словами, разгрузка реактивного уровня по высоким абстракциям не
работает, когда сам реактивный уровень упирается в физический предел
(disturbance не парируется policy, обученной на нормальном трении).
Эксперимент с дополнительным дообучением политики на расширенном
распределении поверхностей частично снимает это ограничение, но
формирует trade-off с пиковой скоростью на отдельно взятой поверхности.

\textbf{Общая трактовка~--- эффект STRL-распределения зависит от среды.}
Полезность каждого верхнего уровня архитектуры не абсолютна, а функционально
определена структурой среды и качеством нижних уровней. В средах со свободным
движением полезность тактического и супервизорного уровней нулевая или
отрицательная (внесение дополнительных степеней свободы без выигрыша
в успехе). В средах с известной геометрией препятствий супервизорный
уровень становится критически важен. В средах с физическими ограничениями
нижнего уровня (низкое трение) ни один из верхних уровней не компенсирует
ограничение.

\textbf{Ограничения работы.} В работе использован простейший
геометрический супервизор без оптимизационного планирования; реактивный
уровень обучен с привилегированным наблюдением (идеальная локализация
и точное знание положения цели); реальная робототехническая платформа
не задействована. Эти ограничения определяют направления развития,
обсуждаемые в заключении.


\section{Заключение}
\label{sec:conclusion}

В работе реализована и эмпирически исследована трёхуровневая иерархическая
архитектура управления мобильным роботом с дифференциальным приводом,
соотносимая с парадигмой STRL школы Осипова-Панова. Реактивный уровень
реализован обученной методом TD3 нейросетевой политикой, тактический~---
LQR-residual с упрощением $e_y = 0$, супервизорный~--- модулем обхода
препятствий по известной карте.

Ablation-анализ на двух семействах сценариев~--- семи конфигураций препятствий
и трёх типах поверхности~--- позволяет количественно охарактеризовать
распределение функций между уровнями. Тактический уровень не меняет
success rate, но улучшает гладкость управления. Супервизорный уровень
разгружает реактивный в задачах с известной геометрией препятствий,
но не справляется в условиях слабого сцепления колёс с поверхностью~---
модельное знание не компенсирует физических ограничений нижнего уровня.
Все результаты воспроизведены при $n = 10$ запусках на детерминированных
сидах со 100\,\% воспроизводимостью.

Направления развития работы естественно следуют из выявленных ограничений.
\textbf{Стратегический уровень} в его полной STRL-формулировке~---
оптимизационный планировщик, выдающий последовательность подцелей на
длинном горизонте, а не supervisorно перехватывающий управление, ---
требует реализации и сравнения с использованным здесь упрощённым
вариантом. \textbf{Visual-language интерфейс} к стратегическому уровню,
в духе прямого предшественника~\cite{veitsenfeld2024}, естественно
дополняет архитектуру визуальной модальностью и связывает её с программой
школы Осипова-Панова в её современной формулировке. \textbf{Реальный
робот} как тестовая платформа~--- ключевое направление, на котором
традиционно сосредоточена работа группы~\cite{kiselev2020dissertation};
перенос архитектуры из PyBullet на платформу с интерфейсом ROS требует
адаптации наблюдения (восстановление позы из SLAM) и проверки устойчивости
обученной политики к sim-to-real разрыву.
\begin{thebibliography}{99}

\bibitem{fujimoto2018td3}
Fujimoto S., van Hoof H., Meger D. Addressing function approximation error in actor-critic methods
// Proc. of ICML. 2018. P.~1587--1596.

\bibitem{mnih2015dqn}
Mnih V., Kavukcuoglu K., Silver D. et al. Human-level control through deep reinforcement learning
// Nature. 2015. V.~518. P.~529--533.

\bibitem{silver2018residual}
Silver T., Allen K., Tenenbaum J., Kaelbling L. Residual policy learning
// arXiv preprint arXiv:1812.06298. 2018.

\bibitem{johannink2019residual}
Johannink T., Bahl S., Nair A. et al. Residual reinforcement learning for robot control
// Proc. of ICRA. 2019. P.~6023--6029.

\bibitem{nachum2018hiro}
Nachum O., Gu S., Lee H., Levine S. Data-efficient hierarchical reinforcement learning
// Proc. of NeurIPS. 2018. P.~3303--3313.

\bibitem{hafner2022director}
Hafner D., Lee K.-H., Fischer I., Abbeel P. Deep hierarchical planning from pixels
// Proc. of NeurIPS. 2022.

\bibitem{sutton1999options}
Sutton R. S., Precup D., Singh S. Between MDPs and semi-MDPs: A framework for temporal abstraction in reinforcement learning
// Artificial Intelligence. 1999. V.~112. P.~181--211.

\bibitem{gat1998three}
Gat E. On three-layer architectures
// Artificial Intelligence and Mobile Robots / D.~Kortenkamp, R.~Bonasso, R.~Murphy (eds.). AAAI Press, 1998. P.~195--210.

\bibitem{kiselev2020dissertation}
Киселёв Г.~А. Интеллектуальная система планирования поведения коалиции робототехнических агентов со STRL архитектурой:
дис. \dots канд. техн. наук. Москва, 2020. 117~с.

\bibitem{kiselev2020}
Киселёв Г.~А., Панов А.~И. Иерархическое психологически-инспирированное планирование для интеллектуальных динамических систем
// Известия РАН. Теория и системы управления. 2020. № 1. С.~123--137.

\bibitem{veitsenfeld2024}
Вейценфельд Д.~А., Киселёв Г.~А. Метод синтеза поведения когнитивного агента на основе обработки мультимодальных сигналов
// Моделирование и анализ данных. 2024. Т.~14. № 4. С.~45--62.

\bibitem{osipov2018sign}
Osipov G. S., Panov A. I. Relationships and operations in a sign-based world model of the actor
// Scientific and Technical Information Processing. 2018. V.~45. P.~317--330.

\bibitem{coumans2021pybullet}
Coumans E., Bai Y. PyBullet, a Python module for physics simulation for games, robotics and machine learning. 2016--2021. URL: \url{http://pybullet.org}

\bibitem{raffin2021sb3}
Raffin A., Hill A., Gleave A. et al. Stable-Baselines3: Reliable reinforcement learning implementations
// Journal of Machine Learning Research. 2021. V.~22. № 268. P.~1--8.

\end{thebibliography}

\end{document}